% =============================================================================
% CVE-2026-82078 (+ CVE-2026-81578) — ARL Case 003 — Xpectra Security Research
% Preamble head: templates/paper-head.tex (verbatim machinery; placeholders
% filled). Every number in this paper comes from designs/facts.json or a cited
% evidence artifact under evidence/ (binding: claims.json at the case root).
% =============================================================================
% !TEX program = pdflatex
\documentclass[journal]{IEEEtran}
\usepackage[utf8]{inputenc}
\usepackage{graphicx}
\usepackage{booktabs}
\usepackage{xcolor}
\usepackage{listings}
\usepackage{hyperref}
\usepackage{microtype}
\usepackage{subcaption}

\hypersetup{colorlinks=true,linkcolor=blue,citecolor=blue,urlcolor=blue}

\definecolor{codebg}{RGB}{245,247,250}
\definecolor{codekw}{RGB}{0,90,150}
\definecolor{codecm}{RGB}{120,120,120}
\definecolor{codestr}{RGB}{160,40,40}

% NOTE: breaklines=true in EVERY style. For pure wire transcripts where a
% single token exceeds the column, add breakatwhitespace=false.
% NOTE: minimal TeX Live has no `yaml` listings language -> yargen style.
\lstdefinestyle{shell}{
  language=bash,
  backgroundcolor=\color{codebg},
  basicstyle=\ttfamily\footnotesize,
  keywordstyle=\color{codekw}\bfseries,
  commentstyle=\color{codecm},
  stringstyle=\color{codestr},
  frame=single,
  breaklines=true,
  columns=fullflexible,
  keepspaces=true,
  showstringspaces=false,
  numbers=left,
  numberstyle=\tiny\color{codecm},
  xleftmargin=1.5em, xrightmargin=1em
}
\lstdefinestyle{smtp}{
  language={},
  backgroundcolor=\color{codebg},
  basicstyle=\ttfamily\footnotesize,
  frame=single,
  breaklines=true,
  breakatwhitespace=false,
  columns=fullflexible,
  keepspaces=true,
  showstringspaces=true
}
\lstdefinestyle{yargen}{
  language={},
  backgroundcolor=\color{codebg},
  basicstyle=\ttfamily\footnotesize,
  frame=single,
  breaklines=true,
  columns=fullflexible,
  keepspaces=true,
  showstringspaces=false
}

\newcommand{\figdir}{../evidence/figures}

\begin{document}

\title{The Chain Dies at Step Three: Measured Three-Arm Closure, Mechanism
Location Without Decompilers, and a Live-Fired Detection Package for the
PaperCut NG Authentication-Bypass to Unsafe-Driver-Loading Chain
(CVE-2026-81578 / CVE-2026-82078)}

% JOURNAL MODE: \IEEEauthorblockN/A are inline dummies in journal mode — stack manually.
\author{Xpectra Security Research\\[0.5ex]
\small Autonomous Red/Blue Operations}

\markboth{The Chain Dies at Step Three: CVE-2026-81578 / CVE-2026-82078}{Xpectra
Security Research \MakeLowercase\expandafter\the\year}

\maketitle

\begin{abstract}
The PaperCut NG/MF chain CVE-2026-81578\,$\rightarrow$\,CVE-2026-82078 was
disclosed by vendor advisory on 2026-08-27, assigned CVEs on 2026-08-28, added
to CISA KEV on 2026-08-31 (due 2026-09-14), and a merged Metasploit module
existed by 2026-09-03. \textbf{This paper makes no discovery claim.} It
contributes what the public record left unmeasured, from a fictional,
zero-egress, three-arm Docker lab built on stock vendor installer bytes:
(i)~an end-to-end reproduction of the full chain on unmodified 25.0.11.75758
using \emph{zero credentials} and a hand-assembled 977-byte proof-of-execution class, with
restore-to-factory configuration proven by cold database reads on repeat runs;
(ii)~a same-day three-arm closure matrix showing the chain fires on 25.0.11 and
dies at the same step (3/11, a 302 bounce of the forged request) on both the
emergency-patch arm 25.0.12-PO-4560.76533 and the maintenance release
25.0.13.76605, plus mechanism location obtained without any decompiler — the
shipped \texttt{papercut.application} service override plus a
\texttt{ValidatingDirectService} class, the vendored Tapestry engine jars being
hash-identical across all arms — and unit corroboration that the payload
artifact still executes on the fixed build, so the fix closes the
\emph{path}, not the artifact;
(iii)~a detection package in which every rule was live-fired or cut, including
an attempt-level wire rule proven to fire on patched arms and a
process-lineage rule that cannot; and (iv)~the non-vacuous-negative
methodology behind these results (cold reads, same-day positive controls, a
five-pillar gate returned \texttt{VALID}), with the 33-entry attempt log
published as the honest-limitations record.
\end{abstract}

\begin{IEEEkeywords}
PaperCut, CVE-2026-82078, CVE-2026-81578, authentication bypass, unsafe
reflection, closure matrix, closed-source fix analysis, detection engineering,
reproducible experiments, honest negative results.
\end{IEEEkeywords}

% =============================================================================
\section{Introduction}
\label{sec:intro}

Print-management servers are attractive targets: they run as a service, hold a
directory of users, and expose a management web interface that must by design
interact with the local system. PaperCut NG/MF is such a product. In late
August 2026 the vendor disclosed a two-link chain: an authentication bypass in
the web management interface (CVE-2026-81578, urgent advisory 2026-08-27, CVE
published 2026-08-28) that an unauthenticated remote request can use to modify
certain system configurations, and an unsafe dynamic class-loading flaw in the
database connection utilities (CVE-2026-82078) which the CVE Record describes
verbatim: the application ``instantiates database driver classes based on
configurable driver names without validating against an allowlist of approved
drivers \ldots\ enables the execution of arbitrary Java bytecode residing on
the application classpath under the security context of the PaperCut server
process''~\cite{cve82078}. Chained, one unauthenticated HTTP session becomes
arbitrary Java bytecode execution as the \texttt{papercut} service user.

The public record is rich. The vendor published mechanism, scores, fixed
versions and per-installer SHA-256 sums~\cite{advisory}; Rapid7 named the
display/execute divergence that glues the chain~\cite{r7}; Huntress reproduced
a pre-auth configuration takeover and the complete RCE on a stock
25.0.11.75758~\cite{huntress}; a full-chain exploit module was merged into
Metasploit on 2026-09-03~\cite{msf}; and both CVEs entered the CISA KEV catalog
on 2026-08-31 with a federal due date of
2026-09-14~\cite{kev}.\footnote{A third-party write-up additionally claims a
bypass of the \emph{first} emergency patch~\cite{techanarchy}; the Metasploit
module header and the Rapid7 post carry the same claim independently. We cite
it, attribute it, and do not adopt it — our lab never measured EPR1, whose
bytes were pulled from the vendor CDN before we could fetch them
(Section~\ref{sec:remediation}).} \textbf{No discovery claim is made anywhere
in this paper}; the credit division is the timeline above.

What the public record does not contain — and what defenders actually need —
is measurement: \emph{which} fixed build kills the chain \emph{at which step};
whether the kill is proven with a live, non-vacuous negative; where the fix
physically lives when the vendor ships no source and no public commit; and what
a defender can actually see on the wire or the host, with false-positive
behavior measured rather than asserted. At intake the public detection surface
for this chain was genuinely empty: no Sigma, Suricata, YARA, ExploitDB, or
vulhub entries, and the official Nuclei template contribution was still an open
pull request~\cite{facts}. That empty surface is now a published, fired package.

\subsection*{Contributions}
\begin{itemize}
\item \textbf{C1 — Measured end-to-end reproduction on stock bytes in a
zero-egress lab.} The full 11-step / 24-request chain fires on an unmodified
vendor installer (25.0.11.75758, our recorded SHA-256 as custody record) with
\emph{zero credentials} (Section~\ref{sec:exploit}). Unlike the merged
Metasploit module — which likewise needs no credentials but ships a
Groovy/HTTP-server payload JAR — our artifact is a 977-byte, hand-assembled
proof-of-execution classfile whose success path throws nothing; it self-deletes its own
dropped bytes, and restore-to-factory \texttt{tbl\_config} is proven by cold
(in-process-lock-free) database reads on three repeat runs and re-proven as a
same-day positive control twice more.
\item \textbf{C2 — A three-arm closure matrix measured in one day, plus
mechanism location without decompilers.} The same unmodified exploit bytes fire
on 25.0.11, die at step 3/11 on emergency-patch arm 25.0.12-PO-4560.76533
(EPR3), and die identically on maintenance release 25.0.13.76605; the 302 death
is captured at raw-packet level (Section~\ref{sec:fix}). The fix is located by
diffing \emph{shipped} artifacts: the stock \texttt{WEB-INF}\slash
\texttt{papercut.application} spec has no direct-service override, both fixed
arms register \texttt{ValidatingDirectService} (byte-identical between arms),
and every vendored Tapestry engine jar is SHA-256-identical across all three
arms — the fix replaces the service class, not the engine. Unit corroboration
inside the fixed arm shows the exact payload bytes still execute when loaded
under the vendor's own JRE: the fix closes the \emph{path}, not the
\emph{artifact} — a transferable method for closed-source fix analysis.
\item \textbf{C3 — A detection package where every shipped rule is live-fired
or cut.} Seven attack SIDs (attempt-, stage-, and execution-labeled) plus a
sensor-health heartbeat SID, six Sigma rules with field aliases, six YARA rules
with clean-stock negatives, and a self-test scanner; three candidate rules were
\emph{cut} with measured reasons, and one attempt rule is proven to fire on
patched arms too (Section~\ref{sec:detection}).
\item \textbf{C4 — The non-vacuous-negative methodology itself}: cold database
reads, same-day positive controls bracketing every negative within minutes, a
five-pillar automated gate returned \texttt{VALID} (10 checks), harness
supersessions logged rather than deleted, and a detector self-test with a
measured non-mutation proof (Section~\ref{sec:lab}).
\end{itemize}

Sections~\ref{sec:char}–\ref{sec:fix} present the vulnerability, lab, and
measurements; Section~\ref{sec:detection} the detection package; Section~\ref{sec:remediation}
patch and IR guidance; Section~\ref{sec:ethics} ethics, honesty, and limits.

% =============================================================================
\section{Vulnerability Characterization}
\label{sec:char}

\subsection{Chain anatomy}
The two CVEs are one kill-chain; neither alone reaches code execution from the
internet. \textbf{Link 1 (CVE-2026-81578):} PaperCut's web tier is Apache
Tapestry~3.0.4 (vendored). A Tapestry \emph{direct service} request encodes a
page and a component listener in the URL. In the affected builds the request
``names one page to render and a different page owning the component to
execute; the access check trusts the \emph{rendered} page''~\cite{r7,facts} —
the shipped public Home/Error pages are unauthenticated, so a request that
displays \texttt{Home} while executing a listener of the privileged
\texttt{ConfigEditor} page runs the listener \emph{before} access validation
completes (the CVE Record's own phrasing: backend actions ``prior to the
completion of access validation checks''~\cite{cve81578}). The reachable effect
is the configuration editor: arbitrary \texttt{tbl\_config} property writes.

\textbf{Link 2 (CVE-2026-82078):} among those writable keys are the external
user-lookup database settings (\texttt{user-lookup.db-driver},
\texttt{db-url}, \texttt{id-to-username-sql}). The connection utility calls
\texttt{Class.forName()} on the configured driver name with no allowlist; the
loaded class's \texttt{<clinit>} runs under the server process's security
context. Because the configuration stage already points the JDBC URL at the
\emph{bundled Derby} engine with an attacker-chosen stored-procedure call
(\texttt{SYSCS\_UTIL.\hspace{0pt}SYSCS\_EXPORT\_QUERY\_\hspace{0pt}LOBS\_TO\_EXTFILE}), the chain also
owns a file-write gadget: it writes a chosen hex BLOB to a CWD-relative path —
\texttt{lib/<name>.class} is on the classpath. The ``arbitrary Java bytecode
residing on the application classpath'' of the CVE wording is thus reachable
without uploading any file through any upload feature.

Vocabulary note, measured rather than assumed: the phrase ``complex direct''
is Rapid7's and the vendor FAQ's, not the product's — a grep for its
spellings across every entry of every extracted jar and the WAR
\texttt{WEB-INF} of all three arms returns \textbf{0 hits in shipped bytes}
(Section~\ref{sec:fix}); the shipped strings that evidence the mechanism are
different (Section~\ref{sec:fixmech}).

\subsection{Trust-hop analysis}
Table~\ref{tab:hops} decomposes the chain by the trust assumption broken at
each hop. Every hop was exercised in our lab (Section~\ref{sec:exploit}); the
chain's length is its detection surface, which is what Section~\ref{sec:detection}
harvests.

\begin{table}[t]
\centering
\caption{Trust hops, each with the measured consequence. ``Assumption broken''
is the implicit trust the build places in something an unauthenticated request
controls.}
\label{tab:hops}
\footnotesize
\begin{tabular}{p{0.225\columnwidth}p{0.32\columnwidth}p{0.28\columnwidth}}
\toprule
Hop & Assumption broken & Measured effect \\
\midrule
H1 anonymous HTTP & management interface is reachable but gated & public
\texttt{Error}/\texttt{Home} pages answer without a session; version oracle
returns \texttt{NG 25.0.11 build 75758} \\
H2 direct-service dispatch & the rendered page authorizes the executed
component & a forged direct-service URL naming \texttt{Home} executes the
privileged listener; response 200 with \texttt{<\,-- Page:\allowbreak Home --\,>} \\
H3 config editor listener & only operators mutate \texttt{tbl\_config} & four
\texttt{user-lookup.*} keys rewritten (driver, URL, SQL, enabled) \\
H4 driver \texttt{Class.\allowbreak forName} & configured drivers are trusted
allow-list members & attacker-named class initializes in the server JVM; a
non-\texttt{Driver} class is tolerated (connect error is caught and logged) \\
H5 Derby procedure & SQL run for lookup is benign & hex BLOB written to
\texttt{lib/<B>.class} (classpath, CWD-relative) \\
\bottomrule
\end{tabular}
\end{table}

\subsection{Access-control analysis of the scores}
The vendor scores the legs separately: 81578 at CVSS:4.0/AV:N/\hspace{0pt}AC:L/\hspace{0pt}AT:N/\hspace{0pt}PR:N/\hspace{0pt}UI:N/\hspace{0pt}VC:L/\hspace{0pt}VI:H/\hspace{0pt}VA:L/\hspace{0pt}SC:N/\hspace{0pt}SI:N/\hspace{0pt}SA:N = 8.8 HIGH and 82078 at \texttt{CVSS:4.0/\hspace{0pt}AV:N/\hspace{0pt}AC:L/\hspace{0pt}AT:N/\hspace{0pt}PR:H/\hspace{0pt}UI:N/\hspace{0pt}VC:H/\hspace{0pt}VI:H/\hspace{0pt}VA:H/\hspace{0pt}SC:H/\hspace{0pt}SI:H/\hspace{0pt}SA:H} = 9.4 CRITICAL~\cite{cve81578,cve82078,advisory}. The \texttt{PR:H} on 82078 encodes ``if an attacker can manipulate system configuration'' — true only via 81578, whose PR:N is what the chain supplies. \textbf{Our measurement matches the chained reading: every leg ran with zero credentials} (Section~\ref{sec:exploit}); no login, no session beyond an anonymous JSESSIONID.

NVD analyzed both CVEs with its own CVSSv3.1 primary metrics: \textbf{81578 =
9.8} (CVSS:3.1/AV:N/\hspace{0pt}AC:L/\hspace{0pt}PR:N/\hspace{0pt}UI:N/\hspace{0pt}S:U/\hspace{0pt}C:H/\hspace{0pt}I:H/\hspace{0pt}A:H) and \textbf{82078 = 9.1} with \texttt{S:C}~\cite{nvd}. These diverge from the vendor v4.0 in effect/scope (81578 \texttt{VC:L}$\to$\texttt{C:H}; 82078 \texttt{S:N}$\to$\texttt{S:C}); we record the divergence, cite both, and blend neither. On our measurements the impact of the chain is ``full CIA inside the service account'': the executed proof-of-execution payload wrote a file as \texttt{uid=1000(papercut)} and ran arbitrary process children as that user — the vendor \texttt{PR:H} understates reachability \emph{for the chain}, which is exactly what the companion CVE is for.

\subsection{Both CWEs, verbatim, as a conflict — no winner}
For CVE-2026-82078 every source agrees: CWE-470, ``Use of Externally-Controlled
Input to Select Classes or Code (`unsafe reflection')''~\cite{cve82078,kev}.
For CVE-2026-81578 the classification \emph{contradicts itself inside the
vendor's own artifacts}, and we present both verbatim rather than invent a
winner:
\begin{itemize}
\item PaperCut's \textbf{CVE Record} (CNA container): ``CWE-305 Authentication
Bypass by Primary Weakness''~\cite{cve81578}. NVD inherits 305 (Secondary) and
adds no CWE of its own~\cite{nvd}.
\item PaperCut's \textbf{own advisory page}: ``CWE-306 Missing Authentication
for Critical Function''~\cite{advisory}. CISA KEV names the entry ``PaperCut
NG/MF \emph{Missing Authentication for Critical Function} Vulnerability'' with
CWE-306~\cite{kev}, and Rapid7 follows 306~\cite{r7}.
\end{itemize}
No source reconciles them~\cite{facts}. Our measurements cannot adjudicate the
taxonomy — both classes describe what we observed (an unauthenticated request
reaching a privileged function; authorization performed \emph{after} dispatch)
— so both stand cited, unlabeled by preference, in this paper.

% =============================================================================
\section{Lab and Methodology}
\label{sec:lab}

\begin{figure}[t]
\centering
\includegraphics[width=\columnwidth]{\figdir/fig-topology.pdf}
\caption{Three-arm lab on one \texttt{--internal} Docker bridge
(\texttt{cve82078\_lab}, 172.30.77.0/24). The network is \texttt{Internal=true}
(counted check), containers are single-homed with no default route, the vendor
installer and first boot run under \texttt{RUN --network=none}, and vendor logs
show zero \texttt{UnknownHost}/\texttt{ConnectException} lines. The only
variable across arms is the installer build; the attacker sits on the bridge
gateway. Evidence reads (\texttt{docker exec}, host-side) are a measurement
channel, not an attack path. Generated figure
(\texttt{evidence/figures/make\_figs.py}).}
\label{fig:topo}
\end{figure}

\subsection{Target provenance}
Table~\ref{tab:prov} records the three installers — the experiment's single
variable. All three were fetched from the vendor CDN and hashed on intake.

\begin{table}[t]
\centering
\caption{Installer provenance (all fetched 2026-09-11; hashes recomputed by
\texttt{claim\_gate}). The stock build's integrity rests on our custody hash
alone — the vendor publishes checksums only for the fixed builds. The first
intake comparison of the patched arm reported MISMATCH because the wrong
bulletin row (MF 76604) was read; re-deriving the NG row programmatically gave
MATCH — a gate that failed correctly and is logged (attempt log \#1).}
\label{tab:prov}
\footnotesize
\begin{tabular}{p{0.185\columnwidth}p{0.148\columnwidth}p{0.295\columnwidth}p{0.16\columnwidth}}
\toprule
Arm (build) & Size (B) & SHA-256 (12-hex head, full in \texttt{ioc/} \S6) & Vendor hash verdict \\
\midrule
stock 25.0.11.75758 (arm A, .20) & 960057398 & \texttt{644957717fcb\ldots} & none published (custody) \\
EPR3 25.0.12-PO-4560.76533 (arm B, .40) & 986902850 & \texttt{68858979b4db\ldots} & MATCH (Wayback row Linux/\allowbreak76533) \\
MR 25.0.13.76605 (arm C, .30) & 986907259 & \texttt{f21311fa3833\ldots} & MATCH (bulletin NG/\allowbreak Linux/\allowbreak76605) \\
\bottomrule
\end{tabular}
\end{table}

\subsection{Lab conditions (S-A set) and integrity controls}
The lab was built once (Dockerfile + per-arm installer), and identical
conditions across arms were the design constraint:
\emph{S-A0 stock integrity:} vendor installer \texttt{--non-interactive} as
user \texttt{papercut}, no config pre-edits; the setup wizard is completed only
by replaying its own Tapestry POST (browser-equivalent, not a bypass shortcut).
\emph{S-A1 database state:} stock \texttt{database.type=Internal} with empty
\texttt{database.driver}; the chain pre-configures nothing — the exploit itself
writes the \texttt{user-lookup.*} rows via the web bypass.
\emph{S-A2 fictional admin:} lab-only admin credentials (free-tier trial
licence, $\leq$5 users, no registration — vendor licence rules
respected~\cite{licence}) exist for the product to be healthy; the exploit never
uses them (zero-credentials is a testable property, not a claim).
\emph{S-A3 isolation:} as Figure~\ref{fig:topo}.
\emph{S-A4 resources:} three arms $\approx$6.2/7.9 GB; fixed arms are stopped
between phases and re-greened before use.

Integrity controls used throughout: \texttt{verify.sh} health gates (25 ok on
stock / 24 ok on non-stock arms, which lack the flag-file check) run
\emph{before and after every} attack run; every evidence file carries
\texttt{\#\#\#\ cmd:}/\texttt{\#\#\#\ at:} headers; and configuration reads are
\emph{COLD} — the app-server is stopped, the Derby table is read
\texttt{ij bootReadOnly}, then the server restarts and re-greens. (The first
two warm \texttt{ij} attempts died on the live JVM's embedded lock and are kept
in the attempt log, \#8, as invalid — a warm read of a locked in-process
database is not evidence.)

\subsection{The agent pipeline}
The case ran as a gated pipeline (intake $\rightarrow$ research $\rightarrow$
lab $\rightarrow$ exploit design/implementation $\rightarrow$ evidence $\rightarrow$
fix verification $\rightarrow$ detection $\rightarrow$ paper), one child agent
per phase under a parent that re-ran gates independently rather than trusting
child self-reports (e.g., the parent re-executed the lab gate to a
\texttt{25 ok/0 fail} green of its own; \texttt{ab\_gate} and
\texttt{claim\_gate} are scripts, not prose). Two rules did real work: (i)~every
state change requires a fresh re-verify before any green is citable — after a
rebuild sequence destroyed a green container mid-run (attempt log \#5), a
self-report of ``green'' would have shipped a phantom; (ii)~invalid attempts
are kept, never rewritten (attempt log \#2 records a confabulated
sink-location line written \emph{before} extraction and later replaced by
byte-level evidence — nothing downstream had consumed it). C4's value is
exactly this machinery made citable; Section~\ref{sec:ethics} publishes the
digest.

% =============================================================================
\section{Exploitation Evidence (C1)}
\label{sec:exploit}

\subsection{The attack, step by step}
\label{sec:pocwalk}

Fig.~\ref{fig:poc} reproduces the chain as it happened on screen, captured
during the live single-take recording (\texttt{videos/takes/t5-show.mov},
SHA-256 \texttt{f9c8b0f8fc5e\ldots}, sealed in the attempt log) and
\emph{never} from the edited film, so no post-production element appears in
any still. The seven times were located by optical character recognition
(OCR) on raw frames before cutting, and the generator
(\texttt{evidence/figures/make\allowbreak\_stills.py}) re-runs that OCR gate
on every final crop, so a caption cannot drift away from its screenshot
(placement log: \texttt{evidence/figures/STILLS-OCR.md}). Step numbers match
the chain anatomy of Fig.~\ref{fig:chain}.

\begin{enumerate}
\item \emph{Oracle} (R1, unauthenticated): one \texttt{curl} to the anonymous
  \texttt{Error} page prints the exact build, \texttt{PaperCut NG 25.0.11
  (Build 75758)}, which is in the KEV-patched range, before any credential exists.
\item \emph{Weapon}: the exploit is a single stdlib script, 1{,}910 lines.
  The excerpt shown is the payload's socket pump, an outbound connect and a
  \texttt{ProcessBuilder} of \texttt{/usr/bin/script -qc /bin/sh -i} (a pty
  that line-buffers the shell), drained with per-direction
  \texttt{available()} polling, the mechanism behind the per-command shell
  responsiveness evidenced in \S\ref{sec:exploit} (kill-shot subsection).
\item \emph{Listener armed} on the gateway (\texttt{172.30.77.1:4444})
  before the chain launches, so the callback cannot race the catcher.
\item \emph{Callback} (T4): the victim calls back from its own address,
  \texttt{CONNECT from 172.30.77.20:50512}, the first outbound proof of code
  execution.
\item \emph{Owned shell}: \texttt{id} answers
  \texttt{uid=1000(papercut)}; \texttt{whoami} answers \texttt{papercut}, and
  \texttt{cat /root/flag.txt} is refused. The shell sits exactly at the
  service account, no privilege bonus.
\item \emph{Defacement re-rendered} (T6/T7): after writing the payload CSS
  and one service restart, the login page is reloaded from a blank tab and
  comes back serving the banner \emph{from the server}, persistent content
  change on the anonymous surface.
\item \emph{Full-control proof}: the flag \texttt{PCUT{[}25.0.11-
  stock-82078{]}} is read (lab evidence read over \texttt{docker exec}, our
  integrity channel, not an in-shell path), and the session transcript is
  SHA-256 sealed before the take ends.
\end{enumerate}

\begin{figure*}[t]
\centering
\begin{subfigure}[b]{0.48\textwidth}
\centering
\includegraphics[width=\textwidth]{\figdir/fig-poc-a.pdf}
\caption{Step 1: version oracle}
\end{subfigure}\hfill
\begin{subfigure}[b]{0.48\textwidth}
\centering
\includegraphics[width=\textwidth]{\figdir/fig-poc-c.pdf}
\caption{Step 3: listener armed}
\end{subfigure}

\vspace{2pt}
\begin{subfigure}[b]{0.48\textwidth}
\centering
\includegraphics[width=\textwidth]{\figdir/fig-poc-d.pdf}
\caption{Step 4: callback from the victim}
\end{subfigure}\hfill
\begin{subfigure}[b]{0.48\textwidth}
\centering
\includegraphics[width=\textwidth]{\figdir/fig-poc-e.pdf}
\caption{Step 5: shell as \texttt{papercut}}
\end{subfigure}

\vspace{2pt}
\begin{subfigure}[b]{0.48\textwidth}
\centering
\includegraphics[width=\textwidth]{\figdir/fig-poc-f.pdf}
\caption{Step 6: defacement served by the victim}
\end{subfigure}\hfill
\begin{subfigure}[b]{0.48\textwidth}
\centering
\includegraphics[width=\textwidth]{\figdir/fig-poc-g.pdf}
\caption{Step 7: full-control flag}
\end{subfigure}
\caption{The attack on screen: six of the seven stills (step 2's code
excerpt is Fig.~\ref{fig:pocb}), cut from the raw recorded take
\texttt{t5-show} with no editing between take and figure. Rows: steps 1/3,
4/5, 6/7. Each frame is OCR-gated against the string quoted in its caption.
Screenshots from our own isolated lab: attacker \texttt{172.30.77.1}, victim
\texttt{172.30.77.20}.}
\label{fig:poc}
\end{figure*}

\begin{figure*}[t]
\centering
\includegraphics[width=\textwidth]{\figdir/fig-poc-b.pdf}
\caption{Step 2: the exploit's reverse-shell pump as it sat in the take,
the payload the agent wrote (SHA-256 \texttt{eb8782e900b7\ldots}) opened at
the relay loop: connect, pty via \texttt{script}, two-way \texttt{available}
pump, single clean exit.}
\label{fig:pocb}
\end{figure*}

\subsection{Wire script}
The exploit is one stdlib-Python script (\texttt{exploit/\allowbreak
papercut\_82078.py},
SHA-256 \texttt{eb8782e900b7\ldots} pinned in this paper), no third-party
dependencies, no attacker-side HTTP service. Its \texttt{--wire-dump} mode
prints the exact requests deterministically (fixed PRNG seed, golden-hashed in
\texttt{exploit/tests/golden/}, 24 request blocks) with no socket opened;
Listing~\ref{lst:wire} shows the two shapes that matter: the anonymous
version/session requests and one forged direct-service write pair. The whole
dry-run chain is 11 logical steps / 24 HTTP requests on one keep-alive
connection. A hard subnet guard (target must be a literal IPv4 in the lab
subnet; no flag bypasses it, hostnames refused) is shipped and refusal-tested
against out-of-lab addresses.

\begin{table*}[t]
\centering
\caption{Closure matrix — same unmodified exploit script (SHA-256 head
\texttt{eb8782e900b7}), same day (2026-09-11, UTC), same payload (vocabulary note: the proof-of-execution PoC artifact appears in raw
captures under its internal capture name);
\texttt{--dry-run} against all three arms. The two fixed arms die at the
\emph{identical} step; the sink leg (82078) is never reached on either.
Artifacts: \texttt{evidence/phase05/run-0[6-9]-*} + \texttt{evidence/pcap/}.}
\label{tab:closure}
\footnotesize
\begin{tabular}{p{0.235\textwidth}p{0.21\textwidth}p{0.24\textwidth}p{0.19\textwidth}}
\toprule
Leg & arm A: stock 25.0.11.75758 & arm B: EPR3 25.0.12-PO-4560.76533 & arm C: MR 25.0.13.76605 \\
\midrule
R1 version oracle (public page) & 200 --- NG 25.0.11/75758 & 200 --- NG
25.0.12/76533 & 200 --- NG 25.0.13/76605 \\
R2 anonymous Home + JSESSIONID & 200 & 200 & 200 \\
\textbf{R3 first forged \texttt{service=direct} POST} & \textbf{200} --- listener
runs un-authed (81578 holds) & \textbf{302} \texttt{Location:\hspace{0pt} /app?\hspace{0pt}service=\hspace{0pt}page/\hspace{0pt}Home} & \textbf{302}, identical \\
R4--T10 (config writes, SYSCS drop, classload) & run & \multicolumn{2}{c}{\emph{never reached}} \\
marker \texttt{/tmp/\allowbreak cve2026-82078-canary} & PRESENT,
\texttt{papercut}:644 & ABSENT (host-side check) & ABSENT (host-side check) \\
COLD \texttt{tbl\_config} pre vs post & restored-to-stock, byte-identical & byte-identical (untouched) & byte-identical (untouched) \\
run.json outcome & \texttt{success=true}, rc 0 & \texttt{error: \ldots(status 302)}, rc 1 & same error, rc 1 \\
\bottomrule
\end{tabular}
\end{table*}

\begin{table*}[t]
\centering
\caption{Latency, in milliseconds, as recorded by the tool itself in
\texttt{run.json} (never hand-copied): \texttt{T1$-$T0} = SYSCS export-ack
(second lookup that runs the export SQL), \texttt{T3$-$T2} = driver-swap
class-load ack — i.e., configuration write to \texttt{<clinit>} complete.
Across all nine measured runs the class-load ack spans 15.2--41.6\,ms; the
export ack (which includes Derby memory-DB boot) spans 89.1--501.3\,ms. The
marker file is observed in-band with the same ack (no out-of-band race assumed).}
\label{tab:latency}
\footnotesize
\begin{tabular}{p{0.29\textwidth}p{0.13\textwidth}p{0.13\textwidth}p{0.175\textwidth}p{0.13\textwidth}}
\toprule
Run & mode & class B & T1$-$T0 ex\-port ack & T3$-$T2 load ack \\
\midrule
run-01 (16:56Z) & dry-run & 977 B & 269.0 & 17.0 \\
run-02 (16:57Z) & dry-run & 977 B & 501.3 & 26.4 \\
run-03 (16:58Z) & dry-run & 977 B & 396.5 & 18.2 \\
run-04c-a (17:12Z) & dry-run & 977 B & 135.6 & 15.2 \\
run-04c-b (17:12Z) & reverse-shell & 1705 B & 121.1 & 27.1 \\
run-07 (17:52Z, positive control) & dry-run & 977 B & 358.1 & 22.1 \\
run-09 (17:57Z, positive control) & dry-run & 977 B & 498.9 & 24.5 \\
run-06-p06-a/b (18:35Z) & dry + shell & 977/1705 B & 496.0 / 89.1 & 27.3 / 41.6 \\
\bottomrule
\end{tabular}
\end{table*}

\begin{table*}[t]
\centering
\caption{Detection master table — every shipped rule, its label class, its
measured fire(s) against real evidence (recorded pcaps/log windows plus fresh
phase-06 live runs), and the named negative control that returned zero.
Attempt-level rules fire on PATCHED arms too (bytes arrive); execution-level
rules require the payload to have run. Transcripts:
\texttt{evidence/rule-fires/}; full ledger \texttt{detection/sigma-fires.md}.}
\label{tab:detect}
\footnotesize
\begin{tabular}{p{0.20\textwidth}p{0.10\textwidth}p{0.325\textwidth}p{0.26\textwidth}}
\toprule
Rule & Class & Fires on (count) & Negative control $\rightarrow$ 0 \\
\midrule
SID 202682041 forged direct POST naming \texttt{user-lookup.} key & attempt &
dry-run pcap $\times$10, rshell $\times$20, fresh live $\times$10/$\times$20;
\textbf{patched negctl $\times$1} & benign GET pcap $\rightarrow$ 0 (health SID
alive) \\
SID 202682042 forged \texttt{\$Form} setter (driver/JDBC values) & stage &
$\times$20 + live & patched pcap has no \texttt{\$Form} setter bytes \\
SID 202682043 SYSCS export SQL config write & stage & $\times$2 + live &
patched arm: zero SYSCS bytes (grep-verified) \\
SID 202682044 \texttt{VALUES CAST(X'CAFEBABE\ldots'} card value & stage &
$\times$2 + live & patched arm: no such bytes; benign 0 \\
SID 202682047 forged direct POST answered \textbf{200} (bypass ACCEPTED) & execution (wire) &
$\times$1 per attack run & both fixed arms $\to$ 0 \\
SID 202682045 SYN server $\rightarrow$ listener ports & execution (wire) &
$\times$1 per reverse-shell run & dry-run \& fixed arms $\rightarrow$ 0 \\
SID 202682046 forged POST answered \textbf{302} page-bounce (BLOCKED) & blue info &
fires exactly on the two fixed-arm negctls $\times$1 & n/a \\
Sigma \texttt{serverlog\_\allowbreak cardid\_\allowbreak hexblob} (server.log sink line) & execution (host) &
24 hits incl. fresh live events & negctl stream (fixed arm) $\to$ 0 \\
Sigma \texttt{derby\_memory\_boot} & execution (host) & 9 hits incl. LIVE-caught boot &
stock \emph{file-dir} boot (no \texttt{memory:}) $\to$ 0 \\
Sigma \texttt{pcapp\_sh\_child} (process lineage kill-shot) & execution (host) &
4 hits (fresh run + run-04c lineage) & 16 benign proc events $\to$ 0 \\
Sigma catch-all (pc-app $\rightarrow$ shell/net tools) & execution (host) & 6 hits &
negctl stream $\rightarrow$ 0 \\
Sigma \texttt{tblconfig\_\allowbreak system\_\allowbreak churn} (+ burst corr.) & slow signal &
12 + 1 correlated alert & negctl stream $\to$ 0 (cannot fire) \\
YARA remnants (6 rules) & post-forensics & attack log/pcap windows incl.\ fresh runs &
\textbf{clean idle} \texttt{server.log}/\texttt{derby.log} windows; benign pcap \\
\bottomrule
\end{tabular}
\end{table*}

\begin{lstlisting}[style=smtp, caption={Wire script, exact golden-dump bytes
(\texttt{exploit/tests/\allowbreak golden/\allowbreak wire-dump-dry-run.bin}; R1, R2, and the first
forged write pair; \texttt{<slot>} is a fresh random per request). The
\texttt{quickFindForm} request selects the config key; the
\texttt{\$Form} request writes the Derby driver name. Percent-decoded
\texttt{\%24} $=$ the \texttt{\$} of Tapestry's form grammar; the wire rule of
Listing~\ref{lst:suricata} keys on the encoding-agnostic substring.},
label={lst:wire}, basicstyle=\ttfamily\scriptsize, xleftmargin=0.9em, xrightmargin=0.3em]
----- WIRE DUMP [R1] (162 bytes, NOT SENT) -----
GET /app?service=page/Error HTTP/1.1
Host: 172.30.77.20:9191
Origin: http://172.30.77.20:9191
Referer: http://172.30.77.20:9191/app
Connection: keep-alive
----- WIRE DUMP [R2.h0] (143 bytes, NOT SENT) -----
GET /app HTTP/1.1                    [ ...headers as above... ]
----- WIRE DUMP [R3.driver.qf] (365 bytes, NOT SENT) -----
POST /app?service=direct/<slot>/Home/ConfigEditor/quickFindForm HTTP/1.1
Content-Type: application/x-www-form-urlencoded
Content-Length: 96

sp=S0&Form0=%24TextField%2CdoQuickFind%2Cclear&%24TextField=user-lookup.db-driver&doQuickFind=Go
----- WIRE DUMP [R3.driver.set] (389 bytes, NOT SENT) -----
POST /app?service=direct/<slot>/Home/ConfigEditor/$Form HTTP/1.1
Content-Length: 127

sp=S1&Form1=%24TextField%240%2C%24Submit%2C%24Submit%240&%24TextField%240=org.apache.derby.jdbc.EmbeddedDriver&%24Submit=Update
\end{lstlisting}

\subsection{Proof-of-execution payload anatomy}
The host carries no \texttt{javac} (JRE-only JDK~25), so the payload is
hand-assembled: a pure-\texttt{struct} ClassFile writer emits a valid Java-8
(major 52) class with explicit constant pool, one \texttt{<clinit>} and one
\texttt{<init>}, and the StackMapTable frames the verifier demands at branch
and handler targets. The dry-run payload is \textbf{977 bytes}; its
\texttt{<clinit>} writes the case marker string to the marker path under
\texttt{/tmp/} via \texttt{Files.\allowbreak write}, deleting its own
\texttt{lib/<B>.class} and \texttt{tmp/<c>.csv} in a \texttt{finally}, and
throws \texttt{RuntimeException} \emph{only} if the write failed — the success
path throws nothing, and the class need not implement
\texttt{java.sql.Driver} (the connect failure after \texttt{Class.forName} is
caught and logged by the application; the response still renders 200, which is
itself the wire signal SID~202682047 keys on). Byte identity under a fixed
class name makes the artifact testable: the golden build is diffable and the
\texttt{--selftest-classfile} mode builds and parses back every shape.

\subsection{Dry-run matrix and restore-to-stock (three runs)}
Three fresh runs against arm A (\texttt{run-01}/\texttt{02}/\texttt{03},
$\approx$70 s apart) each: success \texttt{true}, \texttt{teardown\_errors =
[]}, the proof-of-execution marker file present at
\texttt{/tmp/\allowbreak cve2026-82078-canary} as \texttt{papercut}:644 with the marker
content, and —
after each run — \textbf{COLD} \texttt{tbl\_config} reads byte-identical to the
factory baseline on all seven \texttt{user-lookup.*} keys. Every run uses fresh
random names (class, CSV, memory-DB), so idempotency was proven against a
changing artifact, not one. Post-run \texttt{verify.sh} shows \emph{exactly
one} expected failure (the marker-absence checks are a post-exploit state
the pre-attack gate cannot pass); this honesty detail is kept in the run record rather than
hidden by reordering the gate.

\subsection{Kill-shot: reverse-shell process tree}
The same chain with the connect-back variant (refused by the tool unless a dry
run preceded it in-session) landed a shell inside arm A at 17:12:58Z — 1.47\,s
after the harness listener bound. The payload's own \texttt{ps -ef} (captured
over the shell, and cross-checked by host-side polls of the same window) is the
strongest single artifact of execution:
\begin{lstlisting}[style=shell, caption={Measured process tree (abridged;
\texttt{evidence/processes/run-04c-*-tree.txt} +
\texttt{run-04c-*/reverse-shell-session.txt}, transcript 13017 B). The
application JVM (\texttt{comm} \texttt{pc-app}, Tanuki wrapper +
\texttt{AppServer} in its command line) spawns the interactive shell; the
shell's grandchild is our evidence command. \texttt{uid=1000(papercut)}, CWD
\texttt{server}. A fresh phase-06 re-fire reproduced the identical lineage with
drifted PIDs (2096/1127) — shape constant, PIDs are not.}, label={lst:tree},
language={}, numbers=none]
$ id; pwd
uid=1000(papercut) gid=1000(papercut) ...
/usr/local/papercut/server
  PID  PPID USER     COMMAND
 8694     1 papercut ./app-monitor [conf]
 8696  8694 papercut .../jre/bin/pc-app [flags...]
   ...wrapper.WrapperSimpleApp
        biz.papercut.pcng.server.AppServer
11666  8696 papercut /bin/sh -i  <- payload
11676 11666 papercut ps -ef      <- evidence
\end{lstlisting}

% =============================================================================
\section{Fix Verification (C2): Three-Arm Closure and Mechanism Location}
\label{sec:fix}

\subsection{The closure matrix}
Table~\ref{tab:closure} states the measured result: the chain \emph{fires} on
25.0.11 and \emph{dies at step 3/11} on both fixed arms — the forged
direct-service POST from an unauthenticated session is bounced
(\texttt{302} $\rightarrow$ \texttt{Location: /app?service=page/Home}) before
any config write exists, so the 82078 sink leg is never reached on either arm.
The 302 was verified in the raw-packet capture (server-side tcpdump ascii), not
only in client-side logs, and the POST body carrying the sink payload string
\texttt{user-lookup.db-driver} is present in \emph{both} fixed arms' pcaps —
the bytes arrive; the authorization layer refuses them. This is consistent with
the vendor FAQ's ``EPR3 users are protected'' and the Metasploit header's ``v2
verified to remediate'', and \emph{refutes nothing public}: every published
bypass claim concerned EPR1, not EPR3 or the MR.

\subsection{Mechanism location without decompilers}
\label{sec:fixmech}
No fix commit is public (closed source; the CVE Record's only reference is the
advisory). Per operator ruling we used differential \emph{binary} inspection
only: extraction, hashing, and \texttt{strings} — no decompiler was run, and we
claim nothing about internal code beyond shipped strings. Three artifacts
locate the mechanism:
\begin{enumerate}
\item \textbf{Full jar manifest (345 \texttt{server/lib/*.jar} per arm):}
stock$\rightarrow$fixed moves exactly the vendor's own stems
(\texttt{pcng-server*}, install-helper, log4j 2.25.3$\rightarrow$2.25.5); the
vendored \texttt{tapestry-3.0.4.jar} and \texttt{tapestry-contrib} are
SHA-256-\emph{identical across all three arms} — the fix is \emph{not} in the
engine.
\item \textbf{Per-class hash sets of the web/authz jar set:} both fixed arms
add the same three classes; two are \texttt{ValidatingDirectService} and
\texttt{ValidatingActionService}, byte-identical EPR3\,$=\,$patched. Shipped
strings show the mechanism:
{\scriptsize\texttt{extends org/apache/\allowbreak tapestry/engine/\allowbreak DirectService}}, a
\texttt{validate(IRequestCycle)}
member, \texttt{ownerPage}/\allowbreak\texttt{OWNER\_PAGE\_INDEX} members (it
re-validates the \emph{owner} page of a direct-service invocation), and on the
page base class a rejection string \texttt{`Rejected post-validation access to
page \{\}' from \{\} with IP: \{\}'} plus \texttt{PageRedirectException} — the
302 we measured, from the bytes.
\item \textbf{The shipped application spec (decisive):} the WAR's
\texttt{WEB-INF/papercut.application} is the service-to-class wiring the engine
loads at boot. Stock (\texttt{aa6240c2\ldots}) has \emph{no} override of the
\texttt{direct} service; EPR3 and the maintenance release ship the identical
spec (\texttt{f4bdbf5d\ldots}, byte-identical pair) adding two overrides
(Listing~\ref{lst:spec}); runtime copies inside all three containers
byte-match their shipped specs.
\end{enumerate}
The honest negative travels with the positive: the advisory phrase ``complex
direct'' has \textbf{0 hits} across every entry of every extracted jar and the
WAR specs of all arms — it is third-party vocabulary, and the shipped-string
evidence is Listing~\ref{lst:spec} plus the \texttt{Validating*} classes.
The fixed MR additionally drops three sink-side \texttt{user-lookup} keys in
its external-lookup class (defense-in-depth on leg~2) — EPR3 closes the chain
without it, so we classify that as hardening, not the kill.

\begin{lstlisting}[style=shell, caption={The fix, located by shipped-spec diff
alone (\texttt{evidence/phase05/static/} wp4c): stock adds nothing; EPR3 and
25.0.13 ship the \emph{same} two service overrides; engine jars identical
across arms. This is a mechanism \emph{location} argument from bytes the vendor
ships — not a source-code claim.}, label={lst:spec}]
--- war-stock/WEB-INF/papercut.application
+++ war-epr3/WEB-INF/papercut.application   (= war-patched, byte-identical)
@@ -11,6 +11,8 @@
     <library id="contrib" specification-path="...Contrib.library"/>
+    <service name="direct" class="biz.papercut.pcng.web.
+      services.ValidatingDirectService"/>
+    <service name="action" class="biz.papercut.pcng.web.
+      services.ValidatingActionService"/>
     <service name="custom-content" class="...CustomContentService"/>
\end{lstlisting}

\subsection{Unit corroboration: the fix closes the path, not the artifact}
\label{sec:unitcorr}
A 302 proves the request was refused; it does not prove the sink \emph{would}
have failed. To separate path-closure from artifact-removal we ran the sink's
own call shape \emph{inside the fixed arm}: the exact
\texttt{--generate-classfile} payload bytes (980 B, major 52, SHA-256 head
\texttt{43d60b84\ldots}) were placed in a scratch directory of the EPR3
container and loaded with a 20-line loader calling \texttt{Class.forName(NAME)}
under the arm's \emph{bundled} Corretto-21 JRE. The class initialized:
\texttt{<clinit>} ran, the marker file appeared with content and SHA-256
(\texttt{aeba89c3\ldots}) identical to the stock-arm fire. An ordered negative
control ran first (empty slot + absent class name $\rightarrow$
\texttt{ClassNotFoundException}, slot stays empty), and cleanup was verified.
Two failures precede this result and are kept in the attempt log: the host
JDK-25 \texttt{javac} could not emit the loader (no \texttt{--release} data for
$\leq$21 targets), so the loader was compiled \emph{inside} the container with
the vendor-shipped \texttt{ecj-3.33.0.jar} under the bundled JRE (\#14, both
the failed and successful commands archived); and a first-pass negative control
checked the slot \emph{after} the positive load, proving nothing, and was
superseded by the ordered reproof (\#15). Caveats: the loader ran as root in a
scratch dir (not as the \texttt{papercut} JVM) — owner-\texttt{papercut}
execution is banked from the stock-arm live runs — and the class ``runs its
proof-of-execution body when loaded'' on the fixed build; what the fix removes is the
\emph{load path} the web chain needs. Generalizing: for closed-source targets,
behavioral closure + shipped-artifact diff + in-situ unit corroboration
together locate a fix and characterize what it does \emph{not} change
(``path, not artifact''), without any decompilation.

\subsection{Why the negatives are not vacuous: five pillars}
\texttt{ab\_gate} re-derived the fix-verification evidence from artifacts and
returned \texttt{verdict=VALID checks=10 pillars\_proven=5/5}
(\texttt{evidence/phase05/ab\_gate-output-*.txt}):
\emph{monitor-alive} — both fixed arms answered benign GETs with 200 through
each exploit window, \texttt{gc.log} grew during it, \texttt{verify.sh} green
before and after;
\emph{witnessed-inert} — the sink payload string is present in the
\emph{server-received} pcap bytes of both fixed arms while the marker file is absent
and the config is byte-identical (honesty note: on these arms the chain dies
before the sink's own logger can speak, so inertness is wire-level +
zero-effect; payload-level inertness is instead isolated by
\S\ref{sec:unitcorr});
\emph{unit-corroboration} — \S\ref{sec:unitcorr};
\emph{positive-control} — the stock arm fired the same payload with the same
script \emph{the same day}, bracketing each negative within minutes (EPR3
negative 17:46:42Z $\rightarrow$ stock positive 17:52:34Z; patched negative
17:54:51Z $\rightarrow$ positive 17:57:1xZ): environment drift is excluded by
time, not by assertion;
\emph{config-diff-only-fix} — COLD full \texttt{tbl\_config} and
\texttt{server.properties} across all three arms: every changed key classifies
as install-noise or build-identity (10/959 and 17/959 rows; the 7 extra patched
rows are \texttt{user-lookup.*} keys the MR never seeds), and
\texttt{server.properties} differs by 1/23 keys per pair (the per-install
random bcrypt admin hash) — no functional divergence other than the fix could
explain the behavioral difference.

\begin{figure*}[t]
\centering
\includegraphics[width=0.92\textwidth]{\figdir/fig-chain.pdf}
\caption{The chain as measured, end to end, with the fixed-arm death path
overlaid: steps and response status from the run logs and pcaps
(\texttt{evidence/run-0*}, \texttt{evidence/phase05/run-0*}). The dashed red
band is the 302 measured identically on both fixed arms — everything below it
was never reached there. Generated figure; see also
Table~\ref{tab:closure}.}
\label{fig:chain}
\end{figure*}

% =============================================================================
\section{Detection (C3)}
\label{sec:detection}
Doctrine: a rule that cannot be fired against real evidence is deleted, not
shipped; every row of Table~\ref{tab:detect} carries a fire transcript and a
named negative control. The attack corpus is the phase-04/05 recorded pcaps and
log windows plus three \emph{fresh} live fires produced by the detection phase
itself (reverse-shell run, dry-run, and a live-caught Derby memory-boot
window). Tools: Suricata 8.0.6, YARA 4.5.8, pySigma 1.5.0 (parse/translate
validation: 6 ok / 0 failed) with firing performed by our own stdlib Sigma-semantics
evaluator over 96 events built from real evidence (Sysmon-style + ECS field
aliases replicated), \emph{not} a SIEM engine — a caveat stated on every
transcript header.

\subsection{Attempt versus execution: measured, not claimed}
The fix is authorization-side: the forged bytes still \emph{arrive} on patched
arms (Section~\ref{sec:fix}). So request-content rules are labeled
\texttt{attempt} and \textbf{proven} to fire on a patched arm (SID~202682041
fired $\times$1 against the phase-05 patched negctl pcap), while the
response-pairing rule (forged direct POST answered \texttt{200} $+$
\texttt{<\,-- Page: Home --\,>}) states the stronger \texttt{bypass-accepted}
fact and returned \texttt{0} on both fixed arms, whose 302s instead fire the
informational \texttt{BLOCKED} SID. On the host side the process-creation rule
cannot fire on a patched arm at all (the payload is never loaded). The wire
side honestly cannot claim host execution: the class drop is filesystem, not
wire.

\begin{lstlisting}[style=shell, caption={Full attempt-level wire rule (SID
202682041) exactly as shipped in
\texttt{detection/rules/\allowbreak suricata/\allowbreak cve2026\_82078\_\allowbreak suricata.rules} (8.0.6
\texttt{-T} clean; \texttt{\$PCUT\_PORTS} etc. from the lab yaml — this
Suricata build rejects in-rule \texttt{var} declarations, measured). The
msg's parenthetical carries the honest scope: fires identically on patched
arms.}, label={lst:suricata}, xleftmargin=1.0em, xrightmargin=0.4em]
alert http $EXTERNAL_NET any -> $HOME_NET $PCUT_PORTS ( \
    msg:"CVE-2026-82078 PaperCut attempt - forged service=direct ConfigEditor\
 quickFind POST naming user-lookup config key (attempt: fires on patched\
 arms too)"; \
    flow:to_server,established; \
    content:"POST"; http_method; \
    content:"service=direct/"; http_uri; \
    content:"Home/ConfigEditor"; http_uri; \
    content:"TextField=user-lookup."; http_client_body; \
    flowbits:set,pcut_82078_forged_direct; \
    classtype:attempted-admin; \
    sid:202682041; rev:1; \
    metadata:cve 2026_82078 2026_81578, stage attempt, severity high, \
 confidence high, created 09_11_2026, author Xpectra_Security_Research; )
\end{lstlisting}

\begin{lstlisting}[style=yargen, caption={The process-creation Sigma rule
(abridged to detection logic; full file with aliases, false-positive notes and
tags: \texttt{detection/rules/sigma/\allowbreak cve2026\_82078\_\allowbreak pcapp\_sh\_child.yml};
pySigma-validated). Fired 4$\times$ over real process snapshots incl.\ a fresh
live reverse-shell run; 0 over 16 benign process events. Container debug
(\texttt{docker exec}) parents under entrypoint/init, which axis~1 excludes —
measured in clean-window snapshots, not assumed.}, label={lst:sigma}, xleftmargin=0.7em, xrightmargin=0.3em, basicstyle=\ttfamily\scriptsize]
title: CVE-2026-82078 PaperCut JVM spawns interactive shell (kill-shot)
logsource: {category: process_creation, product: linux}
detection:
  pimg_pcapp_sysmon: {ParentImage|endswith: '/pc-app'}
  pimg_pcapp_ecs:    {parent_process.executable|endswith: '/pc-app'}
  pcmd_wrapper_sysmon:
    ParentCommandLine|contains|all:
      - 'org.tanukisoftware.wrapper'
      - 'papercut'
  pcmd_wrapper_ecs:
    parent_process.command_line|contains|all:
      - 'org.tanukisoftware.wrapper'
      - 'papercut'
  cimg_sh_sysmon:
    Image|endswith: ['/sh', '/bash', '/dash']
  cimg_sh_ecs:
    process.executable|endswith: ['/sh', '/bash', '/dash']
  cargv_i_sysmon: {CommandLine|contains|all: ['sh', '-i']}
  cargv_i_ecs:
    {process.command_line|contains|all: ['sh', '-i']}
  condition: >
    (1 of pimg_pcapp_* or 1 of pcmd_wrapper_*)
    and (1 of cimg_sh_*) and (1 of cargv_i_*)
level: critical
\end{lstlisting}

\subsection{Rules we cut, and why (publishable negatives)}
Three shipped-rule candidates were measured and deleted: (i)~any rule keyed on
\texttt{class load failed} text — the success path \emph{emits none} (measured
across attack windows; the expected sink log line is a caught
\texttt{DatabaseUtils} cardID error, not a load failure); (ii)~a wire rule
claiming the class-file \emph{drop} — at drop time there are no wire bytes (the
CSV lands on disk; the self-deleting class lives seconds), so the filesystem
signal is where that IOC honestly lives (IOC dossier \S1); (iii)~a SYSCS CSV
content rule on the network — the CSV never traverses the wire (verified
against the pcap corpus). Stated gaps we did not paper over: process rules were
evaluated over the harness's own \texttt{ps} snapshots (weaker than kernel
provenance; labeled in each rule), and no wire rule claims host-side execution.

\subsection{Self-test detector with non-mutation proof}
The package ships a stdlib audit tool whose \texttt{--audit} mode is passive
(version threshold + an auth-free direct-shape probe, no form, no payload):
against the stock arm it printed \texttt{VERDICT: VULNERABLE}, and its
\emph{non-mutation} was proven by COLD \texttt{tbl\_config} reads before and
after the audit — byte-identical including \texttt{modified\_date}. Its
\texttt{--active} mode refuses without a consent flag, refuses non-stock lab
targets, and refuses out-of-lab addresses (all refusals transcripted); when
authorized it re-ran the chain to a verified marker file, removed it
(fresh-run rule), and left the lab at 25 ok / 0 fail.

% =============================================================================
\section{Remediation and Incident-Response Guidance}
\label{sec:remediation}

\subsection{Patch levels — including the emergency-patch reality}
\begin{itemize}
\item \textbf{Affected:} ``from 0 before 24.1.10, 25.0.13, 26.0.5'' — the CNA
treats \emph{all} versions as potentially impacted~\cite{cve82078}. (NVD's CPE
ranges disagree, treating the EPR lines and some 25.0.x/26.0.x as not
vulnerable~\cite{nvd}; we record the divergence, cite both, act on neither in
isolation.)
\item \textbf{Emergency-patch history (from the vendor's own updates log,
including Wayback captures of removed rows):} EPR1 2026-08-28 02:10 AEST (NG
25.0.12.76497 / MF 76496, plus 26.0.4 builds); EPR2 the same day for the
v24/v25/v26 interim \texttt{-PO-4560} lines; EPR3 2026-09-01 18:22 AEST,
explicitly ``additional hardening and mitigation against potential attack
chains''; maintenance releases \textbf{24.1.10 / 25.0.13 / 26.0.5} published
\textbf{2026-09-10} and stated to ``replace all emergency patches''.
\item \textbf{The EPR bytes are partly gone.} EPR1 and EPR2 installers now
\texttt{404} on the CDN (our probes; the build table survives only in Wayback
snapshots) — which is why our middle arm is EPR3, and why \emph{we could not}
test the third-party ``EPR1 is bypassable'' claim even though Metasploit's
author tested against it~\cite{msf,techanarchy,r7}. What we measured is that
the closure lands at EPR3: chain fires on 25.0.11, dies at 3/11 on
25.0.12-PO-4560.76533 and on 25.0.13.76605, identically.
\end{itemize}

\textbf{CISO guidance:} treat any build before an emergency patch as exposed;
``patched with EPR1/EPR2'' is \emph{not} a defensible state (bytes withdrawn,
bypass claim carried by three independent sources, KEV deadline 2026-09-14);
\textbf{migrate to a maintenance release (24.1.10 / 25.0.13 / 26.0.5); if you
cannot, be on EPR3 or later} — EPR3 is where we measured the path close.
Verify your installer against the vendor-published SHA-256 (we demonstrated why
an intake gate must read the \emph{right} row: product NG vs MF differ).

\subsection{Using the IOCs honestly}
The IOC dossier tags every indicator \texttt{[LAB]} / \texttt{[GENERIC]} /
\texttt{[VENDOR]} (\texttt{ioc/CVE-2026-82078.ioc.md}). \texttt{[LAB]} items —
our marker path/sha, listener address, exact random class names — identify
\emph{our} runs, not attackers, and must never ship as threat-intel pivots.
Hunt the \texttt{[GENERIC]} shapes: an eight-lowercase \texttt{.class} in
\texttt{server/lib} appearing and vanishing within seconds (fs-watch/EDR, not
scheduled scans); a companion \texttt{tmp/*.csv}; \texttt{Booting Derby \ldots
memory:} where the stock DB is file-directory; the
\texttt{pc-app}$\rightarrow$\texttt{/bin/sh -i} lineage; \texttt{tbl\_config}
\texttt{user-lookup\%} rows churned by \texttt{modified\_by=system} with no
admin session. The vendor's own IOC list (server.log strings, Windows-side
\texttt{data/content/*.cmd/.out} litter, the observed
\texttt{whoami \& ver} $\rightarrow$ SimpleHelp/AnyDesk escalation chain) and
a GreyNoise-reported IP~\cite{nothreat} complete the union — the vendor's
caveat travels verbatim: ``The absence of the above indicators is not
confirmation that a system has not been affected.''

\subsection{IR hunt (measured signal strengths)}
For triage: the server.log sink line is the highest-fidelity single artifact
(our YARA/Sigma fired 24$\times$ on real windows, 0 on clean-stock windows);
the wire pair tells you accepted-vs-blocked at a given instant; the process
lineage is the critical alert (cannot false-fire on patched hosts by
construction); the config-churn row is the slow tail for finding
cleaned-up intrusions after the fast artifacts self-deleted. For
compromise-vs-attempt on exposed \texttt{9191/tcp}: a 302 to
\texttt{service=page/Home} in answer to a \texttt{service=direct} URL after
2026-09-01 is \emph{the fix refusing you} — hunt the paired inbound attempt,
not an incident.

% =============================================================================
\section{Ethics, Reproducibility, and Limitations}
\label{sec:ethics}

\subsection{Ethics}
All exploitation ran in the fictional lab of Section~\ref{sec:lab}: a Docker
\texttt{--internal} network with zero runtime egress, fictional domain and
users, fictional admin credentials the chain never used, no third-party host
touched, and the vendor's free $\leq$5-user licence rules respected. Dates,
scores, and CVE/CWE wording come from the cited primary record
(advisory~2026-08-27; CVEs 2026-08-28; KEV 2026-08-31, due 2026-09-14; MRs
2026-09-10); no discovery claim appears anywhere. The functional PoC ships
per operator ruling with minimal proof-of-execution payloads, a hard subnet guard (no bypass
flag), and an \texttt{--active} mode that refuses out-of-contract targets.
Publication of the pack awaits explicit operator instruction while KEV
exposure persists.

\subsection{Honest limitations, from the attempt log}
The case keeps an append-only attempt log (51 entries); the load-bearing ones:
\begin{itemize}
\item \#1 the intake checksum gate's correct MISMATCH from a wrong bulletin row
(Table~\ref{tab:prov} note); \#2 an early sink-location line written before
extraction was caught by self-audit and replaced by byte evidence;
\item \#5 a rebuild sequence destroyed a green container mid-session — after
which the rule ``no citable green without a fresh re-verify after every state
change'' was enforced, at the cost the paper now reports honestly;
\item \#8 warm \texttt{ij} reads of the locked embedded DB were invalid, forcing
the COLD read procedure this paper depends on;
\item \#12 the first two reverse-shell evidence runs were harness-side capture
failures (listener closed before the payload's output pump; listener crashed at
bind) — the \emph{chain} succeeded both times, the evidence was re-made on a
third run, and the supersessions are logged, not deleted;
\item \#14/\#15 the unit-corroboration loader build and negative-control ordering
failed, were replaced by in-container ecj compilation and an ordered reproof;
\item \#16 a first Derby-boot race capture failed and is shipped marked
\texttt{README-FAILED} (byte-cut windows with stale offsets) — never cited.
\end{itemize}

\subsection{Closed-source and adaptation boundaries}
We never decompiled; every fix statement is a hash, a diff of shipped bytes, or
a shipped string (Section~\ref{sec:fixmech}) — we claim nothing about internal
implementations beyond those artifacts. The reproduction is bound to its
conditions: stock PaperCut \textbf{NG} 25.0.11.75758, Linux x64, free tier, one
/\hspace{0pt}24 lab network; the v26 H2/Groovy branch of the public exploit is
\emph{not} what we reproduced (ours is the v25/24 Derby path), EPR1/EPR2 could
not be tested (bytes withdrawn), and the vendor's own claim that all versions
are affected means arm A's ``pre-patch'' status understates the space of
vulnerable installs we simply did not measure. Single-operator, single-host
RAM/IO constraints (arm parking between phases) do not affect the
chain/closure measurements, which were same-day and bracketed.

\subsection{Reproducibility}
The case pack is self-contained: lab definitions and gates
(\texttt{lab/}), the exploit script with golden wire-hash tests
(\texttt{exploit/}), every run's artifacts with command/time headers
(\texttt{evidence/run-*}, \texttt{evidence/phase05/}), the rule set with fire
transcripts (\texttt{detection/}, \texttt{evidence/rule-fires/}), IOCs, the
attempt log, machine facts (\texttt{designs/facts.json}), and
\texttt{claims.json} binding this paper's claims to those paths (checked by
\texttt{claim\_gate}). Figures are generated from the case data
(\texttt{evidence/figures/make\_figs.py} — box geometry is text-measured and
asserted at build time); video stills are a phase-08 handoff tracked in
\texttt{paper/figures-manifest.md}.

% =============================================================================
\section{Conclusion}
For a KEV entry with real deployments and a merged Metasploit module, the
defensively relevant open questions were measurement questions. We measured
them: the chain works end-to-end on stock bytes with zero credentials
(977-byte hand-assembled classfile; config restore cold-proven); both fixed
arms kill it at the same authorization step, and the kill is located in the
shipped service-override wiring without any decompiler; the payload artifact
still executes on the fixed build, so remediation closes a path, not an
artifact; and a detection package whose every shipped rule fired against real
evidence — including one that honestly fires on patched hosts and three cut
rules whose failure modes we publish. The methodology layer (cold reads,
minute-scale positive bracketing, pillar-gated negatives, an append-only
attempt log) is the reusable contribution: it is what makes each number above
a measurement rather than an assertion.

% =============================================================================
\begin{thebibliography}{99}
\bibitem{advisory} PaperCut, ``Security Bulletin (27 Aug 2026) — Urgent
Security Advisory,'' page last updated 2026-09-10. Captured
\texttt{evidence/raw-research/advisory.txt}.
\bibitem{cve81578} CVE Record CVE-2026-81578 (CNA: PaperCut), published
2026-08-28. Capture: \texttt{evidence/raw-research/cveawg\_81578.json}.
\bibitem{cve82078} CVE Record CVE-2026-82078 (CNA: PaperCut), published
2026-08-28. Capture: \texttt{evidence/raw-research/cveawg\_82078.json}.
\bibitem{nvd} NVD entries for CVE-2026-81578 / CVE-2026-82078 (Analyzed;
CVSSv3.1 primary 9.8 / 9.1). Captures:
\texttt{evidence/raw-research/nvd\_8*.json}.
\bibitem{kev} CISA KEV, catalogVersion 2026.09.10, feed SHA-256 recorded in
\texttt{designs/facts.json}; both CVEs \texttt{dateAdded} 2026-08-31,
\texttt{dueDate} 2026-09-14.
\bibitem{r7} Rapid7, ``PaperCut NG/MF: Critical zero-day exploited in the
wild,'' 2026-08-28. Capture: \texttt{evidence/raw-research/r7.html.txt}.
\bibitem{huntress} Huntress, ``PaperCut actively exploited,'' 2026-08-28.
Capture: \texttt{evidence/raw-research/hunt.txt}.
\bibitem{techanarchy} techanarchy.net, ``From patch to exploit \ldots\ (reverse
engineering a zero-day in PaperCut NG),'' 2026-08-30 — third-party EPR1-bypass
claim, cited and attributed, not adopted.
\bibitem{msf} \texttt{metasploit-framework},
\texttt{papercut\_ng\_\hspace{0pt}external\_\hspace{0pt}user\_\hspace{0pt}lookup\_\hspace{0pt}rce.rb}, merged 2026-09-03
(sfewer-r7), rank Excellent; captured sha256 \texttt{fc6a22a9\ldots} as
\texttt{evidence/raw-research/msf\_module.rb}.
\bibitem{nothreat} GreyNoise reported scanner IP 45.142.193.132 (linked from
the vendor advisory's IOC addendum of 2026-09-10, which itself states the
vendor has not independently verified it).
\bibitem{licence} PaperCut, ``NG Free Licence Questions'' KB (free $\leq$5
users, no registration; last updated 2026-04-13).
\bibitem{facts} ARL Case 003 verified-facts machine record,
\texttt{designs/facts.json} (provenance-keyed; merged 2026-09-11).
\end{thebibliography}

\end{document}
